% !TeX program = xelatex
\documentclass[UTF8,a4paper,10pt,fontset=fandol]{ctexart}
\usepackage[left=18mm,right=18mm,top=20mm,bottom=20mm,headheight=14pt,headsep=7mm]{geometry}
\usepackage{amsmath,amssymb,array,booktabs,longtable,xcolor,enumitem,fancyhdr,lastpage,needspace,fancyvrb}
\usepackage{xurl}
\usepackage[unicode,colorlinks=true,linkcolor=Link,urlcolor=Link,bookmarksnumbered=true]{hyperref}
\definecolor{Navy}{HTML}{203E58}
\definecolor{Link}{HTML}{24678C}
\definecolor{Muted}{HTML}{5E6D78}
\definecolor{Good}{HTML}{2B7053}
\setmainfont{texgyrepagella-regular.otf}[BoldFont=texgyrepagella-bold.otf,
  ItalicFont=texgyrepagella-italic.otf,BoldItalicFont=texgyrepagella-bolditalic.otf]
\setmonofont{lmmono10-regular.otf}[Scale=0.88]
\urlstyle{tt}
\hypersetup{pdftitle={全书独立版：配套清单与重训练核查报告},pdfauthor={NeuralPDE-Opt},
  pdfsubject={书稿章节、代码入口、正式实验和覆盖边界}}
\ctexset{section={format=\Large\bfseries\color{Navy},beforeskip=14pt,afterskip=7pt},subsection={format=\large\bfseries\color{Navy},beforeskip=10pt,afterskip=5pt}}
\setlength{\parindent}{0pt}
\setlength{\parskip}{5pt}
\setlength{\tabcolsep}{3pt}
\setlength{\LTpre}{5pt}
\setlength{\LTpost}{6pt}
\renewcommand{\arraystretch}{1.16}
\linespread{1.13}
\setlength{\emergencystretch}{2em}
\setlist[itemize]{leftmargin=1.7em,itemsep=3pt,topsep=4pt,parsep=0pt}
\pagestyle{fancy}
\fancyhf{}
\fancyhead[L]{\small\color{Muted}NeuralPDE-Opt \quad 书稿配套实验}
\fancyhead[R]{\small\color{Muted}独立版 v1.0.0 核查}
\fancyfoot[L]{\footnotesize\color{Muted}核对版本：2026-10-29}
\fancyfoot[R]{\footnotesize\color{Muted}\thepage\,/\,\pageref*{LastPage}}
\renewcommand{\headrulewidth}{0.3pt}
\renewcommand{\footrulewidth}{0pt}


\begin{document}

{\small\bfseries\color{Link}STANDALONE EDITION / VERIFIED RETRAINING}\par

{\LARGE\bfseries\color{Navy}全书独立版：配套清单}\par

{\LARGE\bfseries\color{Navy}与重训练核查报告}\par

\vspace{5pt}\textcolor{Navy}{\rule{\linewidth}{0.7pt}}\par

核对日期：2026年10月29日（Asia/Shanghai）。本报告对应独立版 v1.0.0 的新工作区和本轮结果；原项目冻结版本继续保留。

\textbf{156组正式神经训练、8组FEM、各章新评测和独立复核全部完成。} 全流程退出码为0，最终3,364项产物身份复验通过；原冻结资产9,550个文件逐项未变。完成度指既定配套范围的交付，不表示全书所有算法已经实现，也不表示各方法均收敛到同一精度。

\section{交付范围与审计结论}

\begingroup\small
\begin{longtable}{>{\raggedright\arraybackslash}p{\dimexpr 0.16\linewidth-2\tabcolsep\relax}>{\raggedright\arraybackslash}p{\dimexpr 0.27\linewidth-2\tabcolsep\relax}>{\raggedright\arraybackslash}p{\dimexpr 0.57\linewidth-2\tabcolsep\relax}}
\toprule
\textbf{部分} & \textbf{本轮正式规模} & \textbf{交付与边界} \\
\midrule
\endfirsthead
\multicolumn{3}{l}{\color{Muted}\footnotesize 续表}\\
\toprule
\textbf{部分} & \textbf{本轮正式规模} & \textbf{交付与边界} \\
\midrule
\endhead
\midrule\multicolumn{3}{r}{\color{Muted}\footnotesize 续下页}\\
\endfoot
\bottomrule
\endlastfoot
第二章 & 18组神经训练 + 8组FEM & PINN、Deep Ritz、范数感知VPINN；1D/2D Poisson。最终迭代模型、独立积分、网格加密和重载核查。 \\[2pt]
第三章 & 57组训练；54个直接预测模型 + 3个残差模型 & 四类算子、POD、物理FNO、VarMiON、HINTS；新固定源项/变源项测试；448次混合求解。 \\[2pt]
第四章基础 & 48组训练；72个阶段模型 & 16条流程×3种子；全部从头训练。新64例问题族测试；在线控制比较统一取前8例。 \\[2pt]
第四章深化 & 12组训练 & 非线性Deep Uzawa、三组稀疏观测FEM-NN；新512个观测位置及独立PDE诊断。 \\[2pt]
第五章 & 21组训练；21个选定模型 & 7条流程×3种子；新128例；控制约束、状态约束分别核查。 \\[2pt]
第1、6章 & 说明性配套 & 沿用正文映射与方法局限，不虚构训练任务。 \\[2pt]
\end{longtable}
\endgroup


\textbf{通过的检查：} 5项封装检查、334项数值回归、Python优化模式检查、第三章21项小预算训练/重载验证，以及各章正式预算、冻结、独立评测、检查点重载和物理指标复算。源码清单317项逐一匹配；本机同步的2,426项证据文件及源码副本均通过SHA256核对。

\textbf{预定预算未缩减：} 第二章每组5000次Adam及名义1000次L-BFGS；第三章基础/POD/VarMiON v1为1600 epochs，物理与HINTS为400 epochs，VarMiON v2为800 epochs。第四、五章保留各自正式配置；第五章问题族算子为300 epochs / 4800 updates，ADMM为30轮外迭代。L-BFGS实际计数可因原有停止准则小于上限，逐次计数保存在训练记录中。

\textbf{固定规则：} 第二章保留最终迭代；第三至五章仅按train/val选择。HINTS的初值来自本工作区新训练的同种子VarMiON。没有使用旧正式权重。测试已查看，后续开发须使用train/val，本轮test只能作为复测。

\clearpage\section{本轮结果与解释}

\subsection{第二章：解误差与有限元基准}

下表相对误差为无量纲数值，非百分数。每项3种子；H1指梯度半范数。

\begingroup\small
\begin{longtable}{>{\raggedright\arraybackslash}p{\dimexpr 0.1\linewidth-2\tabcolsep\relax}>{\raggedright\arraybackslash}p{\dimexpr 0.17\linewidth-2\tabcolsep\relax}>{\raggedright\arraybackslash}p{\dimexpr 0.365\linewidth-2\tabcolsep\relax}>{\raggedright\arraybackslash}p{\dimexpr 0.365\linewidth-2\tabcolsep\relax}}
\toprule
\textbf{维数} & \textbf{方法} & \textbf{相对L2：均值 ± 样本标准差} & \textbf{相对H1：均值 ± 样本标准差} \\
\midrule
\endfirsthead
\multicolumn{4}{l}{\color{Muted}\footnotesize 续表}\\
\toprule
\textbf{维数} & \textbf{方法} & \textbf{相对L2：均值 ± 样本标准差} & \textbf{相对H1：均值 ± 样本标准差} \\
\midrule
\endhead
\midrule\multicolumn{4}{r}{\color{Muted}\footnotesize 续下页}\\
\endfoot
\bottomrule
\endlastfoot
1D & PINN & 1.035e-06 ± 4.82e-07 & 9.558e-06 ± 4.39e-06 \\[2pt]
1D & Deep Ritz & 1.482e-05 ± 1.15e-05 & 1.685e-04 ± 1.24e-04 \\[2pt]
1D & VPINN & 1.247e-05 ± 1.16e-05 & 1.444e-04 ± 1.28e-04 \\[2pt]
2D & PINN & 3.804e-06 ± 1.47e-06 & 1.472e-05 ± 5.25e-06 \\[2pt]
2D & Deep Ritz & 8.491e-05 ± 1.65e-05 & 3.610e-04 ± 7.61e-05 \\[2pt]
2D & VPINN & 7.787e-05 ± 1.28e-05 & 3.968e-04 ± 1.02e-04 \\[2pt]
\end{longtable}
\endgroup


8组FEM展示网格收敛。神经边界最大误差为2.220e-16；加密积分导致的相对L2变化最大4.872e-12。此处神经训练时间与CPU上的FEM装配求解时间口径不同，不能据此作等硬件加速结论。范数感知VPINN的有限测试空间损失也不是连续误差的单独证书。

\subsection{第三章：新测试下的直接算子预测}

主问题为反应系数到解的映射，小扩散系数设置沿用固定配置。固定源项测试128样本；变源项为32个新系数组、每组4源项。下表误差单位为百分比；±为种子间总体标准差，单种子不提供不确定性估计。

\begingroup\small
\begin{longtable}{>{\raggedright\arraybackslash}p{\dimexpr 0.42\linewidth-2\tabcolsep\relax}>{\raggedright\arraybackslash}p{\dimexpr 0.08\linewidth-2\tabcolsep\relax}>{\raggedright\arraybackslash}p{\dimexpr 0.3\linewidth-2\tabcolsep\relax}>{\raggedright\arraybackslash}p{\dimexpr 0.2\linewidth-2\tabcolsep\relax}}
\toprule
\textbf{方法/变体} & \textbf{种子} & \textbf{相对L2（\%）} & \textbf{相对H1（\%）} \\
\midrule
\endfirsthead
\multicolumn{4}{l}{\color{Muted}\footnotesize 续表}\\
\toprule
\textbf{方法/变体} & \textbf{种子} & \textbf{相对L2（\%）} & \textbf{相对H1（\%）} \\
\midrule
\endhead
\midrule\multicolumn{4}{r}{\color{Muted}\footnotesize 续下页}\\
\endfoot
\bottomrule
\endlastfoot
baseline/fno & 3 & 0.4620 ± 0.0273 & 3.2897 \\[2pt]
baseline/deepxde\_fourier4 & 3 & 7.7055 ± 0.5446 & 11.8322 \\[2pt]
baseline/gno & 3 & 1.5940 ± 0.0175 & 21.2400 \\[2pt]
baseline/galerkin & 3 & 0.7062 ± 0.0968 & 3.5672 \\[2pt]
pod/rank\_4 & 3 & 4.6278 ± 0.0070 & 9.3134 \\[2pt]
pod/rank\_8 & 3 & 2.7149 ± 0.0509 & 5.9411 \\[2pt]
pod/rank\_16 & 3 & 2.5273 ± 0.0712 & 5.1668 \\[2pt]
pod/rank\_32 & 3 & 2.6174 ± 0.0390 & 5.2824 \\[2pt]
physics/supervised\_128 & 3 & 0.6849 ± 0.0565 & 5.7625 \\[2pt]
physics/supervised\_32 & 3 & 1.7867 ± 0.1352 & 8.8836 \\[2pt]
physics/energy\_0 & 3 & 0.4519 ± 0.0169 & 1.6847 \\[2pt]
physics/hybrid\_32 & 3 & 0.4727 ± 0.0222 & 1.8211 \\[2pt]
physics/hybrid\_128 & 3 & 0.4372 ± 0.0264 & 1.7720 \\[2pt]
varmion/varmion & 3 & 7.3381 ± 0.0854 & 12.9603 \\[2pt]
varmion/unstructured & 3 & 9.6642 ± 0.2176 & 11.4802 \\[2pt]
varmion\_v2/dense\_data & 3 & 7.2831 ± 0.0575 & 12.9895 \\[2pt]
varmion\_v2/spd\_joint\_aug128 & 3 & 3.2762 ± 0.0858 & 4.9198 \\[2pt]
varmion\_v2/dense\_joint & 1 & 4.1956 ± 0.0000 & 6.6877 \\[2pt]
varmion\_v2/spd\_data & 1 & 6.4915 ± 0.0000 & 11.3480 \\[2pt]
varmion\_v2/spd\_joint & 1 & 5.8365 ± 0.0000 & 7.7016 \\[2pt]
\end{longtable}
\endgroup


FNO与Galerkin本轮L2误差较小，但DeepONet约7.71\%，GNO的H1约21.24\%，原VarMiON约7.34\%；不能用“审计通过”掩盖这些精度差异。SPD联合训练加增广的VarMiON v2为3.28\%，数据、输出空间与预算也与基础模型不同。POD投影和精确约化Galerkin仅为诊断，不计作可部署预测模型。

\subsection{HINTS：离散求解精度与计时范围}

\begingroup\small
\begin{longtable}{>{\raggedright\arraybackslash}p{\dimexpr 0.19\linewidth-2\tabcolsep\relax}>{\raggedright\arraybackslash}p{\dimexpr 0.13\linewidth-2\tabcolsep\relax}>{\raggedright\arraybackslash}p{\dimexpr 0.14\linewidth-2\tabcolsep\relax}>{\raggedright\arraybackslash}p{\dimexpr 0.19\linewidth-2\tabcolsep\relax}>{\raggedright\arraybackslash}p{\dimexpr 0.17\linewidth-2\tabcolsep\relax}>{\raggedright\arraybackslash}p{\dimexpr 0.18\linewidth-2\tabcolsep\relax}}
\toprule
\textbf{模式} & \textbf{收敛数} & \textbf{平均迭代} & \textbf{总时长/例（ms）} & \textbf{对Jacobi比值} & \textbf{最大离散L2} \\
\midrule
\endfirsthead
\multicolumn{6}{l}{\color{Muted}\footnotesize 续表}\\
\toprule
\textbf{模式} & \textbf{收敛数} & \textbf{平均迭代} & \textbf{总时长/例（ms）} & \textbf{对Jacobi比值} & \textbf{最大离散L2} \\
\midrule
\endhead
\midrule\multicolumn{6}{r}{\color{Muted}\footnotesize 续下页}\\
\endfoot
\bottomrule
\endlastfoot
jacobi & 64/64 & 4040.4 & 201.45 & 1.00 & 2.745e-08 \\[2pt]
fixed & 192/192 & 422.4 & 34.95 & 5.76 & 3.356e-09 \\[2pt]
safeguarded & 192/192 & 406.0 & 35.00 & 5.76 & 7.102e-09 \\[2pt]
\end{longtable}
\endgroup


448次方程-求解器组合全部收敛；神经模式各为64源项×3种子，Jacobi仅64例。误差相对同一65×65 P1系统的直接解。连续参考L2平均误差仍约0.09574\%，不能把离散线性系统误差写成连续PDE误差。计时包含装配、条件矩阵、残差转换和接受/回退，每例重复3次取中位数，不含模型加载；约5.76倍只适用于同机Jacobi对照。

\Needspace{8\baselineskip}\subsection{第四章：纯神经、头修正与深化分开报告}

基础交付使用固定17模态线性问题族。下表选择固定算例与新留出前8例的纯神经结果，单位为控制相对L2百分比，±为三种子样本标准差；其余阶段保存在完整证据中。

\begingroup\small
\begin{longtable}{>{\raggedright\arraybackslash}p{\dimexpr 0.34\linewidth-2\tabcolsep\relax}>{\raggedright\arraybackslash}p{\dimexpr 0.33\linewidth-2\tabcolsep\relax}>{\raggedright\arraybackslash}p{\dimexpr 0.33\linewidth-2\tabcolsep\relax}}
\toprule
\textbf{方法（neural阶段）} & \textbf{固定算例误差（\%）} & \textbf{同一8例留出控制误差（\%）} \\
\midrule
\endfirsthead
\multicolumn{3}{l}{\color{Muted}\footnotesize 续表}\\
\toprule
\textbf{方法（neural阶段）} & \textbf{固定算例误差（\%）} & \textbf{同一8例留出控制误差（\%）} \\
\midrule
\endhead
\midrule\multicolumn{3}{r}{\color{Muted}\footnotesize 续下页}\\
\endfoot
\bottomrule
\endlastfoot
direct & 4.89 ± 3.33 & 固定实例方法 \\[2pt]
kkt & 0.3953 ± 0.293 & 固定实例方法 \\[2pt]
alm & 0.5894 ± 0.298 & 固定实例方法 \\[2pt]
fno\_state & 2.174 ± 0.904 & 0.6436 ± 0.0641 \\[2pt]
fno\_kkt & 6.377 ± 4.59 & 1.121 ± 0.116 \\[2pt]
deeponet\_state & 5.71 ± 0.5 & 2.773 ± 0.135 \\[2pt]
deeponet\_kkt & 5.844 ± 1.67 & 5.686 ± 0.444 \\[2pt]
fedonet\_state & 4.157 ± 0.123 & 1.798 ± 0.0872 \\[2pt]
fedonet\_kkt & 2.269 ± 0.398 & 1.968 ± 0.148 \\[2pt]
mimo & 2.489 ± 0.553 & 2.875 ± 0.468 \\[2pt]
rb\_state & 7.538 ± 4.55 & 6.377 ± 0.546 \\[2pt]
rb\_dino & 2.531 ± 1.13 & 1.675 ± 0.211 \\[2pt]
\end{longtable}
\endgroup


输出头修正可在声明的有限模态空间内达到接近舍入误差的控制误差；它是训练后线性头求解的阶段结果，不是纯神经模型自行达到同等精度。基础双层、Deep Uzawa与FEM-NN的固定特征表示也不能等同一般非线性网络。FEM梯度含离散误差，甚至精确谱解的FEM256诊断仍约0.003808。

\textbf{深化实验：} 隐藏层均参与训练。Deep Uzawa采用可训练状态网络与解析控制块；FEM-NN采用可训练源项网络。以下为独立PDE重求后的指标。

\begingroup\small
\begin{longtable}{>{\raggedright\arraybackslash}p{\dimexpr 0.27\linewidth-2\tabcolsep\relax}>{\raggedright\arraybackslash}p{\dimexpr 0.29\linewidth-2\tabcolsep\relax}>{\raggedright\arraybackslash}p{\dimexpr 0.24\linewidth-2\tabcolsep\relax}>{\raggedright\arraybackslash}p{\dimexpr 0.2\linewidth-2\tabcolsep\relax}}
\toprule
\textbf{流程} & \textbf{控制L2（\%）} & \textbf{真实PDE梯度相对量} & \textbf{源项L2（\%）} \\
\midrule
\endfirsthead
\multicolumn{4}{l}{\color{Muted}\footnotesize 续表}\\
\toprule
\textbf{流程} & \textbf{控制L2（\%）} & \textbf{真实PDE梯度相对量} & \textbf{源项L2（\%）} \\
\midrule
\endhead
\midrule\multicolumn{4}{r}{\color{Muted}\footnotesize 续下页}\\
\endfoot
\bottomrule
\endlastfoot
fem\_clean & 1.799 ± 0.041 & 0.0551 & 8.59 \\[2pt]
fem\_noisy & 3.168 ± 0.243 & 0.0369 & 17.32 \\[2pt]
fem\_regularized & 3.117 ± 0.220 & 0.0406 & 13.42 \\[2pt]
uzawa & 1.991 ± 0.067 & 0.0233 & - \\[2pt]
\end{longtable}
\endgroup


稀疏观测1\%噪声时源项误差为17.32\%，加正则后为13.42\%，仍不能仅凭观测拟合宣称辨识正确。512个新观测位置是同一物理问题的空间诊断，不是新问题族泛化；有限内层迭代也不支持套用精确凸内层的收敛保证。

\Needspace{8\baselineskip}\subsection{第五章：约束可行性与最优性残差}

7条流程分成固定实例与问题族。OSNN和ADMM-PINN的独立Gauss控制L2误差分别为0.0049146\%和0.00093698\%；这两个数字不与下表的新问题族测试作横向排名。

\begingroup\small
\begin{longtable}{>{\raggedright\arraybackslash}p{\dimexpr 0.23\linewidth-2\tabcolsep\relax}>{\raggedright\arraybackslash}p{\dimexpr 0.2\linewidth-2\tabcolsep\relax}>{\raggedright\arraybackslash}p{\dimexpr 0.2\linewidth-2\tabcolsep\relax}>{\raggedright\arraybackslash}p{\dimexpr 0.2\linewidth-2\tabcolsep\relax}>{\raggedright\arraybackslash}p{\dimexpr 0.17\linewidth-2\tabcolsep\relax}}
\toprule
\textbf{新128例方法} & \textbf{控制L2（\%）} & \textbf{真实目标相对差（\%）} & \textbf{真实投影残差L2} & \textbf{活动集错分（\%）} \\
\midrule
\endfirsthead
\multicolumn{5}{l}{\color{Muted}\footnotesize 续表}\\
\toprule
\textbf{新128例方法} & \textbf{控制L2（\%）} & \textbf{真实目标相对差（\%）} & \textbf{真实投影残差L2} & \textbf{活动集错分（\%）} \\
\midrule
\endhead
\midrule\multicolumn{5}{r}{\color{Muted}\footnotesize 续下页}\\
\endfoot
\bottomrule
\endlastfoot
Projected KKT-FNO & 4.6641 ± 0.673 & 0.089598 ± 0.0172 & 0.038412 ± 0.00602 & 4.0607 ± 0.732 \\[2pt]
PDOL & 1.6085 ± 0.279 & 0.011903 ± 0.00551 & 0.013344 ± 0.00331 & 1.2885 ± 0.237 \\[2pt]
iUzawa-Net-F & 1.8071 ± 0.134 & 0.015111 ± 0.00288 & 0.013703 ± 0.00112 & 2.1286 ± 0.226 \\[2pt]
iUzawa-Net-S & 3.7659 ± 0.502 & 0.063337 ± 0.019 & 0.029252 ± 0.00451 & 4.106 ± 0.825 \\[2pt]
\end{longtable}
\endgroup


问题族指标采用相同128例、独立真实PDE重求；活动集容差为1e-3。代理固定点残差与真实投影残差分开记录。表中残差非零，说明不能把完成预算解释成严格KKT最优。iUzawa结构中的精确投影和半正定预条件器，也不自动保证网络学习到真实伴随或获得整体收敛定理。

\textbf{CNP是单独的状态上界问题：} 96阶积分目标为0.073962 ± 1.76e-7，FD128状态约束QP参考目标为0.07395772266；对该参考控制的L2误差为1.1473\% ± 0.0928\%。神经状态及真实PDE重求状态的上界违反量均为0。不同离散空间的目标差保留符号；可行性不等于连续最优性。

算子在细网格的控制由固定网格插值/边缘延拓得到，属于独立数值诊断，不是网络分辨率迁移。历史入口名mimo对应本轮Projected KKT-FNO，不能称为原乘法MIMOONet的复刻。

\section{独立运行、选模与保护边界}

独立版可从wheel展开到新的任意工作目录，数值模块不导入原项目或作者仓库。已经从/tmp工作目录核验导入路径。运行仍需要Python、PyTorch、NumPy、SciPy和Matplotlib；完整第三章预算采用CUDA。实测环境为Python 3.12.9、PyTorch 2.5.1+cu121、NumPy 2.0.1、SciPy 1.15.1、Matplotlib 3.10.1、A100 40GB。

本轮独立虚拟环境共享服务器已有科学计算库，未完成全新操作系统或离线全依赖安装验证。DeepXDE LGPL、Galerkin MIT和ocpde\_dl来源声明继续保留；独立不代表全部原创，也不把不同来源改为统一许可证。

正式训练期间没有改写317项数值源码与源清单，也没有降预算或用旧结果补齐。原数据、结果、第二章冻结目录与readable目录共9,550文件、7,868,164,617字节，通过训练前基线逐项比对。源码清单与正式产物清单分别保存，二者不能混作来源许可证证明。

\textbf{恢复协议：} 第二、三章有检查点恢复机制；第四、五章按各自完成审计跳过已完成任务，不能把重新启动某组训练称作精确续训。不得在已有正式工作区重新无条件启动pipeline，也不得手工清除仍有进程持有的运行锁。

\clearpage\section{全书书稿-代码对应清单}

书稿基准为144页PDF；本表页码沿用书稿印刷页，PDF页序号为印刷页码加6。映射来自此前已核对的全书清单：共68条算法/机制映射，不是68个互不重叠的叶节。每条保留覆盖边界；带编号的源码在附录中给出。

“核心覆盖”指主要构造已有实现；“代表实现”限定于声明特例；“部分覆盖/未覆盖”不可当作完整实验；“说明覆盖”不要求另计训练任务。第2至5章已实现并纳入当前正式范围的流程均已完成本轮交付。

\begingroup\small
\begin{longtable}{>{\raggedright\arraybackslash}p{\dimexpr 0.235\linewidth-2\tabcolsep\relax}>{\raggedright\arraybackslash}p{\dimexpr 0.1\linewidth-2\tabcolsep\relax}>{\raggedright\arraybackslash}p{\dimexpr 0.285\linewidth-2\tabcolsep\relax}>{\raggedright\arraybackslash}p{\dimexpr 0.38\linewidth-2\tabcolsep\relax}}
\toprule
\textbf{书稿位置} & \textbf{覆盖层级} & \textbf{本版实现 / 源码索引} & \textbf{边界} \\
\midrule
\endfirsthead
\multicolumn{4}{l}{\color{Muted}\footnotesize 续表}\\
\toprule
\textbf{书稿位置} & \textbf{覆盖层级} & \textbf{本版实现 / 源码索引} & \textbf{边界} \\
\midrule
\endhead
\midrule\multicolumn{4}{r}{\color{Muted}\footnotesize 续下页}\\
\endfoot
\bottomrule
\endlastfoot
1.1：PDE近似与约束优化，1页 & 说明覆盖 & 全书索引连接第2至5章 & 已有四类计算任务的代表配套 \\[2pt]
1.2：经典基准与神经动机，2页 & 说明覆盖 & 第2章FEM、第4章谱/FEM、第5章QP参考 & 有经典对照；不能据单个算例宣称普遍加速 \\[2pt]
1.3：神经方法设置，4页 & 说明覆盖 & 统一符号与逐章导读 \hyperlink{source-1}{[1]} & 函数表示、算子代理、优化与混合角色分别说明 \\[2pt]
1.4：范围与结构，5页 & 说明覆盖 & 本清单与项目当前入口 & 明确四个算法主体和后置分支 \\[2pt]
2.1：引言，7页 & 说明覆盖 & 章节说明 & 给出共同算例和阅读路线 \\[2pt]
2.2.1：共同数值模板，9页 & 代表实现 & mesh.py、fem.py \hyperlink{source-2}{[2]} \hyperlink{source-3}{[3]} & P1 FEM、积分、独立误差与8组网格结果；非所有离散方法 \\[2pt]
2.2.2：可训练表示，12页 & 代表实现 & models.py \hyperlink{source-4}{[4]} & Tanh网络、坐标缩放、边界嵌入；未遍历全部网络表示 \\[2pt]
2.3.1：PINN，14页 & 核心覆盖 & objectives.py \hyperlink{source-5}{[5]} & 6组正式训练；积分强残差平方；硬边界使边界罚项为零 \\[2pt]
2.3.2：修改与训练问题，15页 & 部分覆盖 & models.py、training.py \hyperlink{source-4}{[4]} \hyperlink{source-6}{[6]} & 已有硬边界及Adam/L-BFGS；无损失自适应、残差自适应采样等完整对照 \\[2pt]
2.3.3：误差分析，18页 & 说明覆盖 & metrics.py \hyperlink{source-7}{[7]} & 独立积分、L2/H1半范和物理残差；不是理论误差分解各项的数值证书 \\[2pt]
2.4.1：Deep Ritz，24页 & 核心覆盖 & objectives.py \hyperlink{source-5}{[5]} & 6组；直接最小化能量，不对能量再次平方；1D非齐次边界在仿射可行空间处理 \\[2pt]
2.4.2：弱式/Petrov-Galerkin，26页 & 代表实现 & 同上；mesh.py提供测试基 \hyperlink{source-2}{[2]} & 6组范数感知VPINN，使用Gram矩阵定义弱残差对偶范数；不是WAN \\[2pt]
2.4.3：与PINN比较，29页 & 核心覆盖 & 正式报告 & 同维度共享表示、积分与名义预算；实际调用和成本另记，不宣称等硬件加速 \\[2pt]
2.5.1：域分解、多尺度与混合形式，30页 & 未覆盖 & 当前仅采用全域坐标表示；扩展边界见覆盖核对 & 无XPINN、专门多尺度结构或混合一阶系统完整实验；可选补充 \\[2pt]
2.5.2：架构、正则化与训练动力学，33页 & 部分覆盖 & training.py \hyperlink{source-6}{[6]} & 既定架构、优化历史、种子及恢复检查已有；无系统架构/正则化研究 \\[2pt]
2.5.3：数据增强，34页 & 未覆盖 & 当前源项及边界驱动训练 & 解析解仅作验证；无稀疏/噪声观测同化正式实验 \\[2pt]
2.6：实现与验证，35页 & 核心覆盖 & report.py、修订审查 \hyperlink{source-8}{[8]} & 独立积分、边界、能量一致性、FEM及模型重载已有 \\[2pt]
2.7：优势与局限，37页 & 说明覆盖 & 正式报告解释部分 & 有限测试空间小损失不等于连续误差证书；H1表示梯度半范 \\[2pt]
3.1、3.2.1：解算子与多次查询 & 代表实现 & reaction.py、reaction\_data.py、ReactionOperator \hyperlink{source-9}{[9]} \hyperlink{source-10}{[10]} \hyperlink{source-11}{[11]} & 说明/实验已有：固定扩散的反应系数到解，以及变化源项任务；不覆盖所有目标算子 \\[2pt]
3.2.2：表示、近似与泛化 & 代表实现 & 数据/训练协议、表示误差诊断 \hyperlink{source-12}{[12]} \hyperlink{source-13}{[13]} & train/val/test隔离、归一化、POD表示误差等已有；理论定理不标成“训练完成”，无统一泛化保证 \\[2pt]
3.3.1：DeepONet，式(3.10)的branch-trunk收缩 & 核心覆盖 & DeepXDEOperator及作者核心组件 \hyperlink{source-14}{[14]} \hyperlink{source-15}{[15]} & 已交付，3组；正式ID `deepxde\_fourier4`，trunk输入加入4阶坐标Fourier特征，保留此标注 \\[2pt]
3.3.2：Fourier neural operator & 核心覆盖 & SpectralConv2d / FNO2d \hyperlink{source-16}{[16]} & 已交付，3组；lifting、截断频域卷积、局部分支及projection；不据同网格误差宣称连续算子一致性 \\[2pt]
3.3.2：Graph neural operator，式(3.16) & 核心覆盖 & IntegralKernel2d / GNO2d \hyperlink{source-17}{[17]} & 已交付，3组；17×17图节点、带权消息积分、P1输出重建，梯度误差受输出空间限制 \\[2pt]
3.3.3：Transformer算子；式(3.21)-(3.24)的Galerkin线性注意力 & 核心覆盖 & GalerkinOperator及注意力组件 \hyperlink{source-14}{[14]} \hyperlink{source-18}{[18]} & 已交付，3组；Galerkin注意力＋Fourier解码器。书稿还介绍通用softmax与查询交互，未对每个变体单列当前正式实验 \\[2pt]
3.4.1：物理信息与数据/物理混合目标 & 代表实现 & TorchP1 / objective、训练驱动 \hyperlink{source-19}{[19]} \hyperlink{source-20}{[20]} & 已交付，15组；采用P1能量型物理FNO，含纯数据、无训练标签能量、混合训练。能量方案仍使用val标签选模 \\[2pt]
3.4.2：降阶基/PCA，式(3.28)及固定解码器 & 核心覆盖 & fit\_pod / CoefficientNet / reconstruct \hyperlink{source-21}{[21]} & 已交付，12组；质量内积下的train-only POD，4/8/16/32四个秩。未单列残差贪心选基路线 \\[2pt]
3.4.3：变分结构，式(3.30)-(3.35) & 代表实现 & VariationalMap / UnstructuredMap \hyperlink{source-22}{[22]} & 已交付，6组；固定64维POD基、线性载荷编码、非线性反应系数分支，含非结构化对照 \\[2pt]
3.4.3：变分结构的本项目增强 & 代表实现 & MatrixMap / operator\_loss、trainer\_v2.py \hyperlink{source-23}{[23]} \hyperlink{source-24}{[24]} & 已交付，9组；dense/SPD与算子一致性监督。SPD使用额外物理信息；不是书稿要求的一个独立架构家族，也不是所有变体均3种子 \\[2pt]
3.5.1：学习残差到误差校正，式(3.37)-(3.41) & 核心覆盖 & residual\_data.py、残差模型训练 \hyperlink{source-25}{[25]} \hyperlink{source-26}{[26]} & 已交付，3组残差模型；训练目标来自求解器残差。该节其他组件的缺口见本清单第7节 \\[2pt]
3.5.2：固定HINTS与残差保护，式(3.43)、算法2 & 核心覆盖 & hybrid.solve、test\_hints \hyperlink{source-27}{[27]} \hyperlink{source-13}{[13]} & 已交付，Jacobi/固定融合/保护融合，共448种方程-求解器组合。复用上述3个残差模型，不重复计训练组数 \\[2pt]
3.6：面向目标的实现与验证 & 代表实现 & heldout.py、evaluation.py、audit.py、validation.py \hyperlink{source-28}{[28]} \hyperlink{source-13}{[13]} \hyperlink{source-29}{[29]} \hyperlink{source-30}{[30]} & 当前问题的独立测试、L2/H1半范、物理诊断、模型重载与身份审计已交付；不等于覆盖所有几何/时变/分布外任务 \\[2pt]
3.7：优势与局限 & 说明覆盖 & 阅读导引、原审查报告 \hyperlink{source-31}{[31]} & 说明已有：表示误差、物理信息预算、成本与推广边界，不另设训练任务 \\[2pt]
4.1：引言，73页 & 说明覆盖 & 统一交付 \hyperlink{source-32}{[32]} & 单实例、问题族与混合方法的角色已区分 \\[2pt]
4.2：代表光滑问题，74页 & 代表实现 & problem.py \hyperlink{source-33}{[33]} & 线性二次Poisson控制、低频状态及八阶控制；半线性/抛物型后置 \\[2pt]
4.3：经典数值方法，77页 & 核心覆盖 & fem.py、统一协议 \hyperlink{source-34}{[34]} \hyperlink{source-35}{[35]} & 经典谱优化、P1 FEM和离散伴随；32/64/128/256网格研究，确定性基准不伪造种子 \\[2pt]
4.4.1：direct-PINN，80页 & 代表实现 & instance.py \hyperlink{source-36}{[36]} & 3组；纯神经与头修正阶段分开；硬边界、频率增强及Gauss积分 \\[2pt]
4.4.1：KKT-PINN，80页 & 核心覆盖 & training.py \hyperlink{source-37}{[37]} & 3组；状态、伴随及控制最优性条件训练，独立真实PDE检查 \\[2pt]
4.4.1：增广拉格朗日，80页 & 代表实现 & instance.py \hyperlink{source-36}{[36]} & 3组；外层乘子与神经内层，头修正单列；乘子对应负物理伴随 \\[2pt]
4.4.1：双层神经优化，80页 & 代表实现 & bilevel.py \hyperlink{source-38}{[38]} & 3组；实际训练下层及隐式超梯度，但仍是17个固定正弦特征的线性头基线 \\[2pt]
4.4.2：FNO、DeepONet、FEDONet，83页 & 核心覆盖 & models.py、training.py \hyperlink{source-39}{[39]} \hyperlink{source-40}{[40]} & 三架构各state/KKT角色×3种子＝18组；纯神经和train-only头修正分别记录 \\[2pt]
4.4.2：多输入算子MIMOONet，83页 & 代表实现 & operator\_models.py、operators.py \hyperlink{source-41}{[41]} \hyperlink{source-42}{[42]} & 3组；四个乘法融合branch、共享trunk；模态缩放和17模态输出投影 \\[2pt]
4.4.2：降阶与导数信息RB/DINO，83页 & 代表实现 & 同上operators.py \hyperlink{source-42}{[42]} & RB-state及RB-DINO各3组；DINO有真实导数监督，仍受固定有限空间约束 \\[2pt]
4.4.3：神经策略比较，88页 & 代表实现 & 基础正式报告 & 同一问题族与明确模型阶段；状态代理控制和直接控制另在相同8例上比较 \\[2pt]
4.5.1：Deep Uzawa基线，89页 & 代表实现 & deep\_uzawa.py \hyperlink{source-43}{[43]} & 3组；有限特征基线，保留内层驻点及乘子更新记录 \\[2pt]
4.5.1：Deep Uzawa表示深化，89页 & 代表实现 & 深化训练、深化说明 \hyperlink{source-44}{[44]} & 新3组；全部隐藏层可训练坐标网络，控制块仍为解析二次更新；非精确内层理论复现 \\[2pt]
4.5.2：ISMO，93页 & 核心覆盖 & ismo.py、hybrids.py \hyperlink{source-45}{[45]} \hyperlink{source-46}{[46]} & 3组；非线性代理、真实PDE主动补样后重训；廉价谱PDE不支持一般端到端加速结论 \\[2pt]
4.5.3：FEM-NN基线，96页 & 代表实现 & fem\_nn.py \hyperlink{source-47}{[47]} & 3组；固定特征源项识别、无噪声全场观测，状态由FEM求解 \\[2pt]
4.5.3：FEM-NN观测深化，96页 & 代表实现 & 深化训练、独立评测 \hyperlink{source-44}{[44]} \hyperlink{source-48}{[48]} & 新9组；稀疏无噪声、1\%噪声、相同噪声加正则化各3种子；可训练源项网络和稀疏转置反传 \\[2pt]
4.6：实现与验证，100页 & 核心覆盖 & 基础评测、深化报告 \hyperlink{source-49}{[49]} & 源码/模型冻结、真实PDE重求、目标/梯度、网格/边界及重载审计 \\[2pt]
4.7：讨论与开放问题，102页 & 说明覆盖 & 第四章导读、两份正式报告 \hyperlink{source-50}{[50]} & 表示空间、头修正、有限内层与成本边界均保留 \\[2pt]
5.1：引言，105页 & 说明覆盖 & 统一交付说明 \hyperlink{source-51}{[51]} & 区分控制约束、状态约束、界面和下层变分不等式 \\[2pt]
5.2：非光滑目标，107页 & 部分覆盖 & 问题定义及5.2.1/5.2.2流程 \hyperlink{source-52}{[52]} & 控制盒约束及状态上界有实验；一般非光滑和稀疏控制未全面覆盖 \\[2pt]
5.2.1：OSNN，107页 & 核心覆盖 & models.py、训练 \hyperlink{source-53}{[53]} \hyperlink{source-54}{[54]} & 3组；状态/缩放伴随网络，精确盒投影恢复控制 \\[2pt]
5.2.1：控制分裂ADMM-PINN，107页 & 核心覆盖 & 同上SplitState与训练 \hyperlink{source-54}{[54]} & 3组；保留神经内层、分裂变量和乘子历史；PDE使用内层控制u，可行输出z另报 \\[2pt]
5.2.1：投影KKT-FNO，107页 & 代表实现 & operators.py \hyperlink{source-55}{[55]} & 3组；历史ID为mimo，实际不是第四章乘法分支MIMOONet \\[2pt]
5.2.1：PDOL，107页 & 核心覆盖 & operators.py、reference.py \hyperlink{source-55}{[55]} \hyperlink{source-56}{[56]} & 3组；状态/伴随FNO代理嵌入原始-对偶迭代；代理固定点与真实PDE残差分开 \\[2pt]
5.2.1：iUzawa-Net-F/S，107页 & 核心覆盖 & operator\_blocks.py \hyperlink{source-57}{[57]} & 两种权重展开×3种子＝6组；精确投影、半正定结构预条件器；不保证整个网络收敛 \\[2pt]
5.2.1：L1稀疏控制，107页 & 部分覆盖 & 原模块中的prox工具及测试；见覆盖核对 & 只有基础工具；没有完整稀疏控制训练、参考及独立评测 \\[2pt]
5.2.2：CNP状态约束，112页 & 代表实现 & models.py、独立评测 \hyperlink{source-53}{[53]} \hyperlink{source-58}{[58]} & 3组；光滑可行状态、真实Laplacian恢复控制、独立状态约束QP；不是障碍控制 \\[2pt]
5.2.2：状态分裂ADMM，112页，式(5.19)-(5.21) & 未覆盖 & 尚无完整入口；缺口见覆盖核对 & 后置；控制分裂ADMM和CNP均不能替代此分支 \\[2pt]
5.3.1：界面模型，116页 & 未覆盖 & 尚无正式实现 & 后置；分片PDE、单侧场、值/通量跳跃及独立经典基准未交付 \\[2pt]
5.3.2：界面神经表示，118页 & 未覆盖 & 尚无正式实现 & 后置；DCSNN软约束、lifting/硬嵌入与界面诊断未交付 \\[2pt]
5.4.1：VI与障碍控制，122页 & 未覆盖 & 尚无正式实现 & 后置；下层障碍解、互补条件与独立控制参考未交付 \\[2pt]
5.4.2：神经双层表述，123页 & 未覆盖 & 尚无正式实现 & 后置；不能用第四章光滑有限特征双层或CNP替代 \\[2pt]
5.4.3：Moreau包络与S2-FOBA，124页 & 未覆盖 & 尚无正式实现 & 后置；包络差项、两阶段算法及固定控制状态精修未交付 \\[2pt]
5.5：实现与验证，126页 & 部分覆盖 & protocol.py、assessment.py \hyperlink{source-59}{[59]} \hyperlink{source-58}{[58]} & 已交付5.2当前七条流程的物理/可行性/活动集/重载诊断；不涵盖后置问题 \\[2pt]
5.6：讨论与开放问题，127页 & 说明覆盖 & 正式报告、导读 \hyperlink{source-60}{[60]} & 当前实验与更一般非光滑问题的边界已明确 \\[2pt]
6：总结与展望，129-132页 & 说明覆盖 & 汇总第二至五章方法、验证与局限 & 不另设神经训练；未实现方向不计入完成度 \\[2pt]
\end{longtable}
\endgroup


继续后置：第五章状态分裂ADMM、界面、障碍控制/S2-FOBA，以及半线性/抛物型扩展。神经初值+CG退出主线，Darcy复刻暂停。不能以控制分裂ADMM代替状态分裂ADMM，不能以CNP代替障碍控制。

\Needspace{10\baselineskip}\section{人工复核更正与证据索引}

\subsection{保留原证据，明确两处来源说明}

第四章自动生成的summary.md仍有“复用9组seed0”的历史模板文字。本轮adopted\_jobs.json为空列表，27项原实现任务及21项补充任务均有训练记录，pipeline未传递复用参数。因此本轮48组全部从头训练。该模板句不适用于本轮；原摘要与FINAL\_AUDIT保持原字节身份，另提供带更正的summary.reviewed.md和REVIEW\_NOTES\_ZH.md。本报告引用其数值而不沿用错误来源描述。

INSTALLATION.json记载的是第一次安装的早期wheel与314项源码。正式训练前更新记录\nolinkurl{PRETRAIN_SOURCE_UPDATE.json}、运行依赖核验DEPENDENCIES.json与本轮317项清单共同界定实际训练版本。不能将早期安装hash当作最终wheel身份；本报告下面列明最终值。

\subsection{身份与证据路径}

以下路径相对完整独立工作区R。服务器R为\nolinkurl{/home/yhr/Desktop/NeuralPDE-Opt-independent-v1}。本机的\nolinkurl{standalone/delivery/workspace}是源码与证据镜像，未下载模型权重和数据文件；完整检查点继续保存在服务器。

\begingroup\footnotesize
\begin{longtable}{>{\raggedright\arraybackslash}p{\dimexpr 0.24\linewidth-2\tabcolsep\relax}>{\raggedright\arraybackslash}p{\dimexpr 0.76\linewidth-2\tabcolsep\relax}}
\toprule
\textbf{身份} & \textbf{SHA256} \\
\midrule
\endfirsthead
\multicolumn{2}{l}{\color{Muted}\footnotesize 续表}\\
\toprule
\textbf{身份} & \textbf{SHA256} \\
\midrule
\endhead
\midrule\multicolumn{2}{r}{\color{Muted}\footnotesize 续下页}\\
\endfoot
\bottomrule
\endlastfoot
书稿 book.pdf & \nolinkurl{85191f64ad9f904afa4aed86ce2dcd4bdfc272b2aec21d2c649261c7ef5670e2} \\[2pt]
最终wheel v1.0.0 & \nolinkurl{f7c59bec496898be46b223f4d401b66091c9d65b9b611be0b7b834bdcf9534a9} \\[2pt]
PAYLOAD\_MANIFEST.json & \nolinkurl{e07d6acdeb2a9365952196f996ef8c8c039917f39a219097a8aa847dd77afef6} \\[2pt]
FINAL\_AUDIT.json & \nolinkurl{633c293ff1e5ddfa054555eddab012d75d8d5289fe9a50c465b4430cb5e87598} \\[2pt]
\end{longtable}
\endgroup


\begingroup\footnotesize
\begin{longtable}{>{\raggedright\arraybackslash}p{\dimexpr 0.24\linewidth-2\tabcolsep\relax}>{\raggedright\arraybackslash}p{\dimexpr 0.76\linewidth-2\tabcolsep\relax}}
\toprule
\textbf{内容} & \textbf{R下的相对路径} \\
\midrule
\endfirsthead
\multicolumn{2}{l}{\color{Muted}\footnotesize 续表}\\
\toprule
\textbf{内容} & \textbf{R下的相对路径} \\
\midrule
\endhead
\midrule\multicolumn{2}{r}{\color{Muted}\footnotesize 续下页}\\
\endfoot
\bottomrule
\endlastfoot
全流程与退出码 & \nolinkurl{verification/PIPELINE.json}；\nolinkurl{verification/pipeline_PROCESS.json} \\[2pt]
前置回归/开发验证 & \nolinkurl{verification/CHECKS.json}；\nolinkurl{verification/SMOKE.json} \\[2pt]
正式训练前更新与依赖 & \nolinkurl{verification/PRETRAIN_SOURCE_UPDATE.json}；\nolinkurl{verification/DEPENDENCIES.json} \\[2pt]
全书审计与最终复验 & \nolinkurl{reports/FINAL_AUDIT.json}；\nolinkurl{verification/FINAL_DELIVERY_RECHECK.json} \\[2pt]
原冻结资产保护 & \nolinkurl{verification/ORIGINAL_ASSETS.json}；\nolinkurl{verification/ORIGINAL_ASSETS_FINAL_CHECK.json} \\[2pt]
第二章 & \nolinkurl{results/ch02/full/summary.md}；\nolinkurl{results/ch02/full/DELIVERY_AUDIT.json} \\[2pt]
第三章训练与测试 & \nolinkurl{results/ch03/final_v1_support/training_audit.json}；\nolinkurl{results/ch03/final_test_v1/verification.json} \\[2pt]
第四章基础 & \nolinkurl{results/ch04/unified_delivery_v1/DELIVERY_AUDIT.json}；\nolinkurl{results/ch04/unified_delivery_v1/summary.reviewed.md} \\[2pt]
第四章深化 & \nolinkurl{results/ch04/nonlinear_depth_v1/DELIVERY_AUDIT.json} \\[2pt]
第五章 & \nolinkurl{results/ch05/unified_delivery_v1/DELIVERY_AUDIT.json} \\[2pt]
人工复核附记 & \nolinkurl{reports/REVIEW_NOTES_ZH.md} \\[2pt]
\end{longtable}
\endgroup


\subsection{新工作区复现入口}

以下命令用于另一个全新工作区，不要对当前完成目录重复运行。完整运行使用正式预算；独立目录初始化拒绝覆盖。

\begin{Verbatim}[fontsize=\footnotesize,frame=single]
python -m pip install neuralpde_book-1.0.0-py3-none-any.whl
neuralpde-book init /absolute/path/new-workspace
neuralpde-book check --workspace /absolute/path/new-workspace
neuralpde-book prepare --workspace /absolute/path/new-workspace
neuralpde-book smoke --workspace /absolute/path/new-workspace
neuralpde-book pipeline --workspace /absolute/path/new-workspace
neuralpde-book status --workspace /absolute/path/new-workspace
\end{Verbatim}

\Needspace{10\baselineskip}\section{源码索引}

路径均相对R，编号对应全书映射表。此索引引用已训练的317项冻结源码清单，未为排版修改数值实现。

\begingroup\footnotesize
\begin{longtable}{>{\raggedright\arraybackslash}p{\dimexpr 0.08\linewidth-2\tabcolsep\relax}>{\raggedright\arraybackslash}p{\dimexpr 0.92\linewidth-2\tabcolsep\relax}}
\toprule
\textbf{编号} & \textbf{源码相对路径} \\
\midrule
\endfirsthead
\multicolumn{2}{l}{\color{Muted}\footnotesize 续表}\\
\toprule
\textbf{编号} & \textbf{源码相对路径} \\
\midrule
\endhead
\midrule\multicolumn{2}{r}{\color{Muted}\footnotesize 续下页}\\
\endfoot
\bottomrule
\endlastfoot
\leavevmode\hypertarget{source-1}{[1]} & \nolinkurl{README.md} \\[2pt]
\leavevmode\hypertarget{source-2}{[2]} & \nolinkurl{src/neuralpde_ch02/mesh.py} \\[2pt]
\leavevmode\hypertarget{source-3}{[3]} & \nolinkurl{src/neuralpde_ch02/fem.py} \\[2pt]
\leavevmode\hypertarget{source-4}{[4]} & \nolinkurl{src/neuralpde_ch02/models.py} \\[2pt]
\leavevmode\hypertarget{source-5}{[5]} & \nolinkurl{src/neuralpde_ch02/objectives.py} \\[2pt]
\leavevmode\hypertarget{source-6}{[6]} & \nolinkurl{src/neuralpde_ch02/training.py} \\[2pt]
\leavevmode\hypertarget{source-7}{[7]} & \nolinkurl{src/neuralpde_ch02/metrics.py} \\[2pt]
\leavevmode\hypertarget{source-8}{[8]} & \nolinkurl{src/neuralpde_ch02/report.py} \\[2pt]
\leavevmode\hypertarget{source-9}{[9]} & \nolinkurl{src/neuralpde_opt/chapter3/reaction.py} \\[2pt]
\leavevmode\hypertarget{source-10}{[10]} & \nolinkurl{src/neuralpde_opt/chapter3/reaction_data.py} \\[2pt]
\leavevmode\hypertarget{source-11}{[11]} & \nolinkurl{src/neuralpde_opt/chapter3/reaction_operators.py} \\[2pt]
\leavevmode\hypertarget{source-12}{[12]} & \nolinkurl{src/neuralpde_opt/chapter3/reaction_training.py} \\[2pt]
\leavevmode\hypertarget{source-13}{[13]} & \nolinkurl{chapters/ch03/evaluation.py} \\[2pt]
\leavevmode\hypertarget{source-14}{[14]} & \nolinkurl{src/neuralpde_opt/chapter3/upstream_operators.py} \\[2pt]
\leavevmode\hypertarget{source-15}{[15]} & \nolinkurl{src/neuralpde_opt/chapter3/upstream_vendor/deepxde_core.py} \\[2pt]
\leavevmode\hypertarget{source-16}{[16]} & \nolinkurl{src/neuralpde_opt/chapter3/fno.py} \\[2pt]
\leavevmode\hypertarget{source-17}{[17]} & \nolinkurl{src/neuralpde_opt/chapter3/gno.py} \\[2pt]
\leavevmode\hypertarget{source-18}{[18]} & \nolinkurl{src/neuralpde_opt/chapter3/upstream_vendor/galerkin_layers.py} \\[2pt]
\leavevmode\hypertarget{source-19}{[19]} & \nolinkurl{experiments/ch03_physics_v1/physics.py} \\[2pt]
\leavevmode\hypertarget{source-20}{[20]} & \nolinkurl{experiments/ch03_physics_v1/run.py} \\[2pt]
\leavevmode\hypertarget{source-21}{[21]} & \nolinkurl{experiments/ch03_pod_v1/pod.py} \\[2pt]
\leavevmode\hypertarget{source-22}{[22]} & \nolinkurl{experiments/ch03_varmion_v1/model.py} \\[2pt]
\leavevmode\hypertarget{source-23}{[23]} & \nolinkurl{experiments/ch03_varmion_v2/model_v2.py} \\[2pt]
\leavevmode\hypertarget{source-24}{[24]} & \nolinkurl{experiments/ch03_varmion_v2/trainer_v2.py} \\[2pt]
\leavevmode\hypertarget{source-25}{[25]} & \nolinkurl{experiments/ch03_hints_v1/residual_data.py} \\[2pt]
\leavevmode\hypertarget{source-26}{[26]} & \nolinkurl{experiments/ch03_hints_v1/run.py} \\[2pt]
\leavevmode\hypertarget{source-27}{[27]} & \nolinkurl{chapters/ch03/hybrid.py} \\[2pt]
\leavevmode\hypertarget{source-28}{[28]} & \nolinkurl{chapters/ch03/heldout.py} \\[2pt]
\leavevmode\hypertarget{source-29}{[29]} & \nolinkurl{chapters/ch03/audit.py} \\[2pt]
\leavevmode\hypertarget{source-30}{[30]} & \nolinkurl{chapters/ch03/validation.py} \\[2pt]
\leavevmode\hypertarget{source-31}{[31]} & \nolinkurl{docs/CH03_READING_ZH.md} \\[2pt]
\leavevmode\hypertarget{source-32}{[32]} & \nolinkurl{chapters/ch04_delivery/README.md} \\[2pt]
\leavevmode\hypertarget{source-33}{[33]} & \nolinkurl{chapters/ch04/problem.py} \\[2pt]
\leavevmode\hypertarget{source-34}{[34]} & \nolinkurl{chapters/ch04/fem.py} \\[2pt]
\leavevmode\hypertarget{source-35}{[35]} & \nolinkurl{chapters/ch04_delivery/protocol.py} \\[2pt]
\leavevmode\hypertarget{source-36}{[36]} & \nolinkurl{chapters/ch04_adaptation/instance.py} \\[2pt]
\leavevmode\hypertarget{source-37}{[37]} & \nolinkurl{chapters/ch04/training.py} \\[2pt]
\leavevmode\hypertarget{source-38}{[38]} & \nolinkurl{chapters/ch04_completion/bilevel.py} \\[2pt]
\leavevmode\hypertarget{source-39}{[39]} & \nolinkurl{chapters/ch04_operator_adaptation/models.py} \\[2pt]
\leavevmode\hypertarget{source-40}{[40]} & \nolinkurl{chapters/ch04_operator_adaptation/training.py} \\[2pt]
\leavevmode\hypertarget{source-41}{[41]} & \nolinkurl{chapters/ch04_adaptation/operator_models.py} \\[2pt]
\leavevmode\hypertarget{source-42}{[42]} & \nolinkurl{chapters/ch04_delivery/operators.py} \\[2pt]
\leavevmode\hypertarget{source-43}{[43]} & \nolinkurl{chapters/ch04_completion/deep_uzawa.py} \\[2pt]
\leavevmode\hypertarget{source-44}{[44]} & \nolinkurl{chapters/ch04_depth/training.py} \\[2pt]
\leavevmode\hypertarget{source-45}{[45]} & \nolinkurl{chapters/ch04_completion/ismo.py} \\[2pt]
\leavevmode\hypertarget{source-46}{[46]} & \nolinkurl{chapters/ch04_delivery/hybrids.py} \\[2pt]
\leavevmode\hypertarget{source-47}{[47]} & \nolinkurl{chapters/ch04_completion/fem_nn.py} \\[2pt]
\leavevmode\hypertarget{source-48}{[48]} & \nolinkurl{chapters/ch04_depth/assessment.py} \\[2pt]
\leavevmode\hypertarget{source-49}{[49]} & \nolinkurl{chapters/ch04_delivery/assessment.py} \\[2pt]
\leavevmode\hypertarget{source-50}{[50]} & \nolinkurl{docs/CH04_READING_ZH.md} \\[2pt]
\leavevmode\hypertarget{source-51}{[51]} & \nolinkurl{chapters/ch05_delivery/README.md} \\[2pt]
\leavevmode\hypertarget{source-52}{[52]} & \nolinkurl{chapters/ch05/problem.py} \\[2pt]
\leavevmode\hypertarget{source-53}{[53]} & \nolinkurl{chapters/ch05/models.py} \\[2pt]
\leavevmode\hypertarget{source-54}{[54]} & \nolinkurl{chapters/ch05_delivery/training.py} \\[2pt]
\leavevmode\hypertarget{source-55}{[55]} & \nolinkurl{chapters/ch05/operators.py} \\[2pt]
\leavevmode\hypertarget{source-56}{[56]} & \nolinkurl{chapters/ch05/reference.py} \\[2pt]
\leavevmode\hypertarget{source-57}{[57]} & \nolinkurl{chapters/ch05/operator_blocks.py} \\[2pt]
\leavevmode\hypertarget{source-58}{[58]} & \nolinkurl{chapters/ch05_delivery/assessment.py} \\[2pt]
\leavevmode\hypertarget{source-59}{[59]} & \nolinkurl{chapters/ch05_delivery/protocol.py} \\[2pt]
\leavevmode\hypertarget{source-60}{[60]} & \nolinkurl{docs/CH05_READING_ZH.md} \\[2pt]
\end{longtable}
\endgroup


\end{document}
